\documentclass[11pt,answers]{exam}
\usepackage[spanish]{babel}
\usepackage[utf8]{inputenc}
\usepackage[T1]{fontenc}
\usepackage{amsmath,amssymb,amsfonts}
\usepackage{graphicx}
\usepackage{colortbl}
\usepackage{xcolor}
\usepackage{multirow}
\usepackage{float}
\usepackage{enumitem}
\usepackage{algorithm}
\usepackage{mathrsfs}
\usepackage{array} % Para controlar el ancho de las celdas
\usepackage{enumitem} % Para personalizar listas
\usepackage{listings}% http://ctan.org/pkg/listings
\usepackage{hyperref} % for hyperlinks
\usepackage{amsmath} % para matemáticas mejoradas
\usepackage{algorithm} % para escribir pseudocódigos
\usepackage{algpseudocode} % para escribir pseudocódigos
\usepackage[version=4]{mhchem}
\usepackage{stmaryrd}

% \usepackage{etoolbox}

% \AtBeginEnvironment{align}{\setcounter{equation}{0}}

\renewcommand{\solutiontitle}{\noindent\textbf{Solución:}\par\noindent}

\lstset{
  basicstyle=\ttfamily,
  mathescape
}

\graphicspath{{public/}}

\setlength{\topmargin}{-.5in} \setlength{\textheight}{9.25in}
\setlength{\oddsidemargin}{-0.5in} \setlength{\textwidth}{7.2in}


\begin{document}
\begin{center}
    \newcommand{\HRule}{\rule{\linewidth}{0.5mm}}
    \begin{minipage}{0.48\textwidth} 
        \begin{flushleft}
            \includegraphics[scale = 0.08]{../../../public/logo_unam.png}
        \end{flushleft}
    \end{minipage}
    \begin{minipage}{0.48\textwidth} 
        \begin{flushright}
            \includegraphics[scale =0.22]{../../../public/logo_ciencias.png}
        \end{flushright}
    \end{minipage}
    \vspace*{-1.5cm}						
    \textsc{\huge Nacional Autónoma de México \\ \vspace{-4px} Universidad }\\[2cm]	
    \textsc{\LARGE Facultad de Ciencias}\\[1.5cm]
    \vspace*{1cm}					
        \HRule \\[0.7cm]							
            { \huge \bfseries Tarea 03}\\[0.4cm]	
        \HRule \\[1.5cm]						    
    \begin{minipage}{0.52\textwidth}													
        \begin{flushleft} \large	
            \small
            \vspace{-0.6cm}	
            \vspace{-0.6cm}	
                \emph{Alumno:}\\
               Ramírez López Alvaro. 316276355\\
            \vspace*{2cm}
        \end{flushleft}																		
        \end{minipage}		
    \begin{minipage}{0.46\textwidth}		
        \vspace{-0.6cm}											
        \begin{flushright} \large						
            \small										
            \emph{Profesor:} Jesús Villagómez Chávez	\\
            \emph{Ayudantes:}
                Gabriela Peña Franco	 \\
                Martha Rubí Gutiérrez González	 \\
        \end{flushright}																
    \end{minipage}	
    \vspace*{1cm}
    \vspace{2cm}
    \begin{center}						
        {\large 9 de septiembre de 2024}
    \end{center}  						
\end{center}	
\textbf{}
\newpage

\begin{enumerate}

    \item ¿Cuáles de las siguientes relaciones son funciones? En caso de ser función, calcula su dominio y su imagen:
    \begin{enumerate}
        \item $\{(n,m) \in \mathbb{Z} \times \mathbb{Z} : n, m \geq 0 \land 5n = m\}$.
        \item $\{(n,m) \in \mathbb{Z} \times \mathbb{Z} : n, m \geq 0 \land 5m = n\}$.
        \item $\{(n,m) \in \mathbb{Z} \times \mathbb{Z} : n, m \geq 0 \land m \leq n\}$.
        \item $\{(n,m) \in \mathbb{Z} \times \mathbb{Z} : n \geq 0 \land m = 3\}$.
        \item $\{(n,m) \in \mathbb{Z} \times \mathbb{Z} : m = n^2\}$.
        \item $\{(n,m) \in \mathbb{Z} \times \mathbb{Z} : m^2 = n^2\}$.
        \item $\{(n,m) \in \mathbb{Z} \times \mathbb{Z} : 4n + 2m = 6\}$.
    \end{enumerate}

    \item Determina la inyectividad, suprayectividad y biyectividad de las siguientes funciones:
    \begin{enumerate}
        \item $f : \mathbb{N} \rightarrow \mathbb{N}, f(n) = 2n$.
        \item $f : \mathbb{N} \rightarrow \mathbb{N}, f(n) = n + 7$.
        \item $f : \mathbb{Z} \rightarrow \mathbb{Z}, f(n) = n + 7$.
        \item $f : A \rightarrow A/R, f(a) = [a]_R$, donde $A$ es un conjunto y $R$ una relación de equivalencia sobre $A$.
    \end{enumerate}

    \item Sea $f : A \rightarrow B$ una función. Demuestra que:
    \begin{enumerate}
        \item $f$ es inyectiva si y sólo si $f^{-1}[f[X]] = X$, para todo $X \subseteq A$.
        \item $f$ es inyectiva si y sólo si $f[X \cap Y] = f[X] \cap f[Y]$, para $X, Y \subseteq A$.
        \item $f$ es suprayectiva si y sólo si $f[f^{-1}[Y]] = Y$, para todo $Y \subseteq B$.
        \item $f$ es biyectiva si y sólo si $f[X^c] = (f[X])^c$, para todo $X \subseteq A$.
    \end{enumerate}

    \item Responde las siguientes preguntas:
    \begin{enumerate}
        \item ¿Existe $g : \mathbb{N} \rightarrow \mathbb{N}$ función tal que $g \neq \text{Id}_{\mathbb{N}}$ y $g \circ g = g$?
        \item ¿Existe $g : \mathbb{N} \rightarrow \mathbb{N}$ función biyectiva tal que $g \neq \text{Id}_{\mathbb{N}}$ y $g \circ g = g$?
        \item ¿Existe $g : \mathbb{N} \rightarrow \mathbb{N}$ función biyectiva tal que $g \neq \text{Id}_{\mathbb{N}}$ y $g \circ g = \text{Id}_{\mathbb{N}}$?
    \end{enumerate}

    \item (Extra) Sea $f : A \rightarrow B$ una función. Definimos la asignación $F : B \rightarrow A$ con regla de correspondencia $F(Y) = f^{-1}[Y]$. Demuestra que:
    \begin{enumerate}
        \item $F$ es función.
        \item Si $f$ es inyectiva, entonces $F$ es suprayectiva.
        \item Si $f$ es suprayectiva, entonces $F$ es inyectiva.
        \item Si $F$ es suprayectiva, entonces $f$ es inyectiva.
        \item Si $F$ es inyectiva, entonces $F$ es suprayectiva.
    \end{enumerate}

\end{enumerate}

\end{document}